\chapter{Image solénoïdale et résolution de Navier--Stokes}

\section{Domaine hydrodynamique et publication de Galerkin}

On fixe
\[
 H^s_\sigma(\mathbb T^3)=
 \{u\in H^s(\mathbb T^3;\mathbb R^3):
   \nabla\!\cdot u=0,\ \int_{\mathbb T^3}u=0\},
 \qquad s>\frac52,\quad\nu>0,
\]
avec
$u_0\in C^\infty\cap H^s_\sigma$ et
$f\in C^\infty([0,\infty);H^\infty_\sigma)
\cap L^2_{\rm loc}([0,\infty);H^{s-1}_\sigma)$. Pour la projection de
Leray--Galerkin $P_M$, on pose
\[
 A_M=-P_M\Delta,
 \qquad B_M(v,w)=P_MP_\sigma((v\cdot\nabla)w),
 \qquad N_M(u):=B_M(u,u),
\]
et le champ fini satisfait
\begin{equation}\label{eq:pdf2-ns-galerkin}
 \partial_tu_M+\nu A_Mu_M+B_M(u_M,u_M)=f_M,
 \qquad u_M(0)=P_Mu_0.
\end{equation}
L'équation \eqref{eq:pdf2-ns-galerkin} est une publication de l'état, non sa
définition complète.

Pour que la notation ne change pas au milieu de la chaîne, on fixe une seule fois les
quatre fonctionnelles d'énergie d'ordre $s$ :
\begin{equation}\label{eq:pdf2-ns-energy-functionals}
 \begin{aligned}
 \mathscr E_{s,M}(u)&:=\frac12\|\Lambda^su\|_{L^2}^2,
 &\mathscr D_{s,M}(u)&:=\|A_M^{1/2}\Lambda^su\|_{L^2}^2,\\
 \mathscr B_{s,M}(u)&:=
   \operatorname{Re}\langle\Lambda^sN_M(u),\Lambda^su\rangle,
 &\mathscr F_{s,M}(u,f)&:=
   \operatorname{Re}\langle\Lambda^sP_Mf,\Lambda^su\rangle.
 \end{aligned}
\end{equation}
L'indice supplémentaire $j$ désigne les mêmes quatre fonctionnelles évaluées sur
la coordonnée publiée du niveau de profondeur $j$ ; il n'introduit pas d'autres
fonctionnelles.
L'équation de Galerkin implique exactement
\begin{equation}\label{eq:pdf2-ns-energy-identity}
 \frac d{dt}\mathscr E_{s,M}(u_M)
 +\nu\mathscr D_{s,M}(u_M)+\mathscr B_{s,M}(u_M)
 =\mathscr F_{s,M}(u_M,f).
\end{equation}

\section{Objet source et coordonnées effectivement construites}

L'image solénoïdale ne s'ajoute pas a posteriori. C'est exactement la construction
autonome $\mathfrak G_{\rm sol}$ produite avant ce chapitre et publiée dans
\eqref{eq:pdf2-g-sol-explicito} ; elle conserve l'état
de Galerkin, les applications de raffinement, les secteurs orientés, l'opérateur dissipatif,
l'advection projetée, la retenue, la counité $E_9$, l'inclusion $J_9$, les
projections $P_9=J_9E_9$ et $Q_{\rm micro}=I-P_9$, la mémoire, le chemin et le terminal.
En particulier,
\[
 E_9J_9=I,
 \qquad P_9^2=P_9,
 \qquad Q_{\rm micro}=I-P_9.
\]
Ces identités typent la décomposition entre publication macroscopique et mémoire
microscopique ; elles ne constituent pas à elles seules une borne globale.

À une profondeur nonadique $j$, l'état conserve
\begin{equation}\label{eq:pdf2-ns-estado}
 \Xi_{M,j}^{\rm NS}=
 \bigl(u_M;q_{A,M,j},q_{V,M,j},T_{M,j};
 \begin{gathered}
 \text{triades, feuillet, orientation, retenue, frontière,}\\
 \text{route, archive et mémoire}
 \end{gathered}
 \bigr).
\end{equation}
Le retour enrichi réunit toutes les branches avant l'évaluation :
\begin{equation}\label{eq:pdf2-ns-retorno}
 \mathcal R_{M,J,t}^{\rm enr}
 =\mathcal R_{1,M}\mathfrak U_{M,J,t}^*,
 \qquad
 \mathcal R_{M,J,t}^{\rm enr}Z_{M,J,t}=u_M(t).
\end{equation}
Une branche isolée de norme $3^J$ n'est pas le retour.

Pour $N>M$, le défaut de commutation avec la non-linéarité est la retenue
\begin{equation}\label{eq:pdf2-ns-carry}
 C_{M\leftarrow N}:=
 P_MB(u_N,u_N)-B_M(P_Mu_N,P_Mu_N).
\end{equation}
La poser égale à zéro sans prouver la commutation élimine précisément l'interaction des
hautes fréquences que le raffinement doit conserver.

Les deux mémoires positives restent séparées. Avec
$r=\pi_{\rm HMT}/\varphi_{\rm HMT}^2>1$,
\begin{equation}\label{eq:pdf2-ns-dh}
 D_M=(r-1)(\mathcal M_M^+-\mathcal M_M^-),
 \qquad
 H_M=(1+r)(\mathcal M_M^++\mathcal M_M^-),
\end{equation}
et
\[
 \dot D_M=\mathscr B_{s,M},
 \qquad
 \dot H_M=\frac{1+r}{r-1}|\mathscr B_{s,M}|.
\]
$D_M$ publie le signe ; $H_M$ conserve la grandeur positive. Les fusionner détruit
l'orientation ou la positivité.

\section{Tour PC--Fock complète et stockage terminal}

L'élévation n'est pas tronquée au degré deux :
\begin{equation}\label{eq:pdf2-ns-pcf}
 \mathcal V_M^{\rm PCF}(u)
 :=\bigoplus_{n\ge0}\frac{R_{\rm F}^n}{\sqrt{n!}}u^{\odot n},
 \qquad
 \|\mathcal V_M^{\rm PCF}(u)\|^2
 =e^{R_{\rm F}^2\|u\|_{H^s}^2}.
\end{equation}
Le rayon fait partie de la coordonnée et ne peut pas disparaître du lecteur. Pour
tout $\rho>0$, on fixe
\begin{equation}\label{eq:pdf2-ns-normalized-degree-one-reader}
 \rho^{[\rho]}_1:=\rho^{-1}\pi_1:
 \Gamma_s^\rho(V)\longrightarrow V,
 \qquad
 \rho^{[\rho]}_1\mathcal V_\rho^{\rm PCF}(u)=u.
\end{equation}
Le niveau initial utilise $\rho^{[R_{\rm F}]}_1$ et le niveau $j$ utilise
$\rho^{[r_j]}_1$ ; cette normalisation est maintenue lors du retrait de la profondeur.
Sur la diagonale cohérente, $z_{n,M}=u_M^{\odot n}$ satisfait
\begin{equation}\label{eq:pdf2-ns-carleman}
 \dot z_{n,M}=L_{n,M}z_{n,M}
 +\mathcal N_{n,n+1;M}z_{n+1,M}
 +\mathcal F_{n,n-1;M}(f_M)z_{n-1,M}.
\end{equation}
Les lecteurs de degrés un et deux sont
\[
 z_M(u)=\nu^{1/2}A_M^{1/2}\Lambda^su,
 \qquad
 w_M(u)=\nu^{-1/2}A_M^{-1/2}\Lambda^sN_M(u),
\]
et Young donne
\begin{equation}\label{eq:pdf2-ns-young-port}
 |\mathscr B_{s,M}(u)|\le
 \beta\nu\mathscr D_{s,M}(u)+\frac{\|w_M(u)\|^2}{4\beta}.
\end{equation}
L'estimation uniforme des lignes de Fourier contrôle ce lecteur quadratique dans
l'état initial. Elle ne contrôle pas encore ses propagations futures.

Soit
\[
 \mathcal N_{M,j}:=
 \frac{r-1}{r+1}H_{M,j}-D_{M,j}^{\rm mem}
 =2(r-1)M_{M,j}^-\ge0.
\]
Sur un horizon $T$, le stockage terminal est
\begin{equation}\label{eq:pdf2-ns-storage}
 \begin{aligned}
 G_{M,j,T}(t)
 &=\mathscr E_{s,M,j}(t)+2(r-1)(M^-_{M,j}(T)-M^-_{M,j}(t)),\\
 \dot G_{M,j,T}+\nu\mathscr D_{s,M,j}+|\mathscr B_{s,M,j}|
 &=\mathscr F_{s,M,j}.
 \end{aligned}
\end{equation}
Pour typer le terminal, on conserve dans l'état solénoïdal complet la coordonnée
de feuillet
\[
 q^-_{M,j}:=\sqrt{2(r-1)\dot M^-_{M,j}}
 \in L^2(0,T;\mathbb R).
\]
Soit $\mathscr X^{\rm sol}_{M,j,T}$ l'espace d'histoire enrichie qui contient,
comme termes orthogonaux, la trace $2^{-1/2}\Lambda^su_{M,j}$ et $q^-_{M,j}$.
La colonne terminale est l'opérateur de coordonnées
\begin{equation}\label{eq:pdf2-ns-terminal-column}
 \begin{aligned}
 \mathcal E_{M,j,T}^{\rm term}(t):\mathscr X^{\rm sol}_{M,j,T}
 &\longrightarrow L^2_\sigma(\mathbb T^3)
   \oplus L^2(t,T;\mathbb R),\\
 \mathcal E_{M,j,T}^{\rm term}(t)\Xi
 &:=2^{-1/2}\Lambda^su_{M,j}(t)
 \oplus\mathbf1_{[t,T]}q^-_{M,j}.
 \end{aligned}
\end{equation}
Par construction des deux coordonnées orthogonales,
\begin{equation}\label{eq:pdf2-ns-terminal-column-gram}
 \|\mathcal E_{M,j,T}^{\rm term}(t)\Xi\|^2
 =G_{M,j,T}(t).
\end{equation}
Le terme de mémoire future est une sortie terminale typée, non une composante de la
racine initiale ; son coût ne sera contrôlé qu'après le télescopage de la colonne
complète. Pour une observation $\tau$, les lignes de perte vivent dans
$[0,\tau]$ et cette mémoire terminale dans $(\tau,T]$ ; les supports temporels sont
disjoints. Dans l'estimation intégrale, on prend $\tau=T$, et alors le terme
futur est nul.

\section{Construction source--first avec relèvement mixte état--force}

L'entrée est fixée avant toute coupure, profondeur et histoire :
\begin{equation}\label{eq:pdf2-ns-data-space}
 \mathcal D_T:=H^s_\sigma(\mathbb T^3)
 \times L^2(0,T;H^{s-1}_\sigma(\mathbb T^3)),
 \qquad d=(u_0,f).
\end{equation}
Comme les champs sont de moyenne nulle, $A=-\Delta$ est positif et inversible dans le
secteur solénoïdal. On introduit l'effort, le flux visqueux et les deux ondes
\begin{equation}\label{eq:pdf2-ns-force-waves}
 \begin{aligned}
 a_T(f)&:=\frac{1}{2\sqrt\nu}A^{-1/2}\Lambda^sP_\sigma f,
 &z_M(u)&:=\sqrt\nu A_M^{1/2}\Lambda^su,\\
 b_M(u,f)&:=z_M(u)-P_Ma_T(f),\\
 q_M^{B,+}(u)&:=\bigl(\mathscr B_{s,M}(u)\bigr)_+^{1/2},
 &q_M^{B,-}(u)&:=\bigl(-\mathscr B_{s,M}(u)\bigr)_+^{1/2},\\
 q_M^B(u)&:=\bigl(q_M^{B,+}(u),q_M^{B,-}(u)\bigr).
 \end{aligned}
\end{equation}
Tous appartiennent à $L^2$ sur l'intervalle correspondant et
$\|q_M^B(u)\|_{\mathbb C^2}^2=|\mathscr B_{s,M}(u)|$.
La puissance croisée est typée par
\begin{equation}\label{eq:pdf2-ns-cross-power}
 \mathscr F_{s,M}=2\operatorname{Re}\langle z_M,P_Ma_T(f)\rangle,
 \qquad \nu\mathscr D_{s,M}=\|z_M\|^2.
\end{equation}
C'est donc \eqref{eq:pdf2-ns-storage}, et non une substitution de la puissance par la
norme de la force, qui donne l'identité exacte de diffusion
\begin{equation}\label{eq:pdf2-ns-scattering-balance}
 \dot G_{M,j,T}+\|b_M\|^2+\|q^B_M\|^2=\|P_Ma_T(f)\|^2.
\end{equation}

\subsection{Signature chronologique, opérateur régularisant uniforme et création mixte}

La donnée de force vit dans $F:=H^{s-1}_\sigma(\mathbb T^3)$, tandis que la
tour PC--Fock utilise comme espace à une particule
$V:=H^s_\sigma(\mathbb T^3)$. L'identification entre les deux n'est pas laissée à
l'inclusion de Galerkin, dont la norme croîtrait avec $M$. Dans le secteur de moyenne nulle, on
fixe l'isomorphisme d'ordre moins un
\begin{equation}\label{eq:pdf2-ns-uniform-smoother}
 \begin{aligned}
 \mathscr S&:=A^{-1/2}:F\longrightarrow V,\\
 \mathscr S_M&:=A_M^{-1/2}P_M=P_M\mathscr S,\\
 \sup_M\bigl(\|\mathscr S_M\|_{F\to V}
 +\|\mathscr S_M^{-1}\|_{P_MV\to P_MF}\bigr)&<\infty.
 \end{aligned}
\end{equation}
La dernière borne est l'équivalence uniforme des normes de Sobolev dans le
sous-espace sans mode zéro. Ainsi, $g:=\mathscr Sf\in L^2(0,T;V)$ et on n'utilise pas
l'inclusion non uniforme $P_MH^{s-1}\hookrightarrow P_MH^s$.

Soit $\Delta_k(T)=\{0<t_1<\cdots<t_k<T\}$. L'espace de signature de la force
régularisée est
\begin{equation}\label{eq:pdf2-ns-signature-space}
 \mathfrak S_T(V):=\bigoplus_{k\geq0}L^2(\Delta_k(T);V^{\widehat\otimes k}),
 \qquad \Delta_0(T)=\{*\}.
\end{equation}
Pour $g=\mathscr Sf\in L^2(0,T;V)$, on définit, sans utiliser la solution,
\begin{equation}\label{eq:pdf2-ns-force-signature}
 \operatorname{Sig}_T(g)_0=1,
 \qquad
 \operatorname{Sig}_T(g)_k(t_1,\ldots,t_k)
 =g(t_k)\widehat\otimes\cdots\widehat\otimes g(t_1).
\end{equation}
La symétrie de l'intégrande scalaire fournit
\begin{equation}\label{eq:pdf2-ns-signature-norm}
 \|\operatorname{Sig}_T(g)\|^2
 =\sum_{k\geq0}\frac{\|g\|_{L^2(0,T;V)}^{2k}}{k!}
 =e^{\|g\|_{L^2(0,T;V)}^2}
 \le e^{c_{\mathscr S}^2\|f\|_{L^2(0,T;F)}^2}.
\end{equation}

Le terme forcé de \eqref{eq:pdf2-ns-carleman} n'est pas un port additif dans
$f$ uniquement. Pour $0<r_{\rm F}<R_{\rm F}$, il est typé au moyen de la création mixte
\begin{equation}\label{eq:pdf2-ns-mixed-creation}
 \begin{aligned}
 \widetilde{\mathcal F}^{R_{\rm F},r_{\rm F}}_M:
 P_MV\widehat\otimes\Gamma_s^{R_{\rm F}}(P_MV)
 &\longrightarrow\Gamma_s^{r_{\rm F}}(P_MV),\\
 \widetilde{\mathcal F}^{R_{\rm F},r_{\rm F}}_M(g\otimes x)
 &:=a^\dagger(g)x.
 \end{aligned}
\end{equation}
Au degré $n$, son action est précisément
$g\otimes z_{n-1}\mapsto n g\odot z_{n-1}$. On conserve ainsi le facteur mixte
$\mathscr Sf\otimes x$ que ne peut produire une transition linéaire sur la somme directe
``état $\oplus$ force''.
Avec les normes pondérées de \eqref{eq:pdf2-ns-pcf},
\begin{equation}\label{eq:pdf2-ns-creation-uniform}
 \|\widetilde{\mathcal F}^{R,r}_M(g\otimes x)\|_{\Gamma^r}
 \le c(R,r)\|g\|_V\|x\|_{\Gamma^R},
 \qquad
 c(R,r):=\sup_{n\ge1}\sqrt n\,r\left(\frac rR\right)^{n-1}<\infty.
\end{equation}
Pour itérer, on fixe $0<\theta<1$ et $r_j=R_{\rm F}\theta^j$ ; chaque pas utilise
$(R,r)=(r_j,r_{j+1})$. Alors
$\sup_{j,M}c(r_j,r_{j+1})\le
R_{\rm F}\sup_{n\ge1}\sqrt n\,\theta^n<\infty$.

La tour s'écrit dans la coordonnée conjuguée
$v_M=\mathscr S_Mu_M$. Soit
$\mathscr S_M^{\Gamma}:=\bigoplus_{n\ge0}
\mathscr S_M^{\widehat\otimes n}$. Le générateur homogène conjugué est
\begin{equation}\label{eq:pdf2-ns-conjugated-generator}
 \mathcal A^{0,\mathscr S}_{M}
 :=\mathscr S_M^{\Gamma}\mathcal A^0_M
   (\mathscr S_M^{\Gamma})^{-1},
 \qquad
 (\mathcal A^{0,\mathscr S}_{M})_{n,m}
 =\mathscr S_M^{\widehat\otimes n}
  (\mathcal A^0_M)_{n,m}
  (\mathscr S_M^{-1})^{\widehat\otimes m}.
\end{equation}
Cette conjugaison est source--first et n'altère pas l'ODE : le lecteur de degré un
applique $\mathscr S_M^{-1}$ à la fin.

Soit $\mathcal T_M^{0,\mathscr S}(t,s)$ le propagateur de la tour conjuguée, inclus
dans la réalisation solénoïdale complète. Pour $R>r>0$, il est d'abord défini sur
\begin{equation}\label{eq:pdf2-ns-radial-transfer}
 \delta_{R\to r}(x_0,x_1,x_2,\ldots)
 :=\bigl(x_0,(r/R)x_1,(r/R)^2x_2,\ldots\bigr),
 \qquad
 \delta_{R\to r}\mathcal V_R^{\rm PCF}(u)=\mathcal V_r^{\rm PCF}(u).
\end{equation}
Cet opérateur est contractif et fixe le vide. Sur
\begin{equation}\label{eq:pdf2-ns-dyson-domain}
 \Gamma_s^R(P_MV)_{\rm alg}\widehat\otimes
 \mathfrak S_T(P_MV)_{\rm fin}
\end{equation}
Soit $\mathfrak S_{T,{\rm fin}}^{\rm dec}$ l'espace engendré par les noyaux élémentaires
\[
 g_k(t_1,\ldots,t_k)=
 g_k^{(k)}(t_k)\widehat\otimes\cdots\widehat\otimes g_k^{(1)}(t_1).
\]
Il est dense dans chaque terme $L^2(\Delta_k(T);V^{\widehat\otimes k})$. Dans la
formule suivante, les symboles $g_k^{(i)}$ sont lus dans chaque noyau élémentaire et
l'action s'étend linéairement à leur espace engendré ; on ne factorise donc pas
arbitrairement un tenseur intriqué. Le lecteur chronologique prend ses valeurs dans
$\Gamma_s^r(P_MV)$ et est défini par
\begin{equation}\label{eq:pdf2-ns-dyson-reader}
 \begin{aligned}
 \mathscr D_{M,k}^{R\to r}(t)(x\otimes g_k)
 :={}&\delta_{R\to r}\int_{\Delta_k(t)}
 \mathcal T_M^{0,\mathscr S}(t,t_k)
 \widetilde{\mathcal F}_M(P_Mg_k^{(k)}(t_k)\otimes\cdot)\cdots\\
 &\hspace{18mm}\cdots\mathcal T_M^{0,\mathscr S}(t_2,t_1)
 \widetilde{\mathcal F}_M(P_Mg_k^{(1)}(t_1)\otimes\cdot)
 \mathcal T_M^{0,\mathscr S}(t_1,0)x\,d\mathbf t,\\
 \mathscr D_{M,T}^{R\to r}(t)&:=\sum_{k\geq0}
 \mathscr D_{M,k}^{R\to r}(t).
 \end{aligned}
\end{equation}
En répartissant la perte totale $R-r$ entre les $k$ créations et les propagateurs,
l'estimation analytique d'échelle de
\eqref{eq:pdf2-ns-creation-uniform} donne, pour des constantes finies dépendant de la
coupure mais non du noyau élémentaire,
\begin{equation}\label{eq:pdf2-ns-dyson-homogeneous-bound}
 \|\mathscr D_{M,k}^{R\to r}(t)(x\otimes g_k)\|_{\Gamma^r}
 \le A_{M,R,r,T}\,
 \frac{B_{M,R,r,T}^{k}}{\sqrt{\Gamma(k+1)}}
 \|x\|_{\Gamma^R}\|g_k\|_{L^2(\Delta_k;V^{\widehat\otimes k})}.
\end{equation}
Le facteur $1/\Gamma(k+1)$ du volume du simplexe et Cauchy--Schwarz dans la
somme des degrés donnent
\begin{equation}\label{eq:pdf2-ns-dyson-bound}
 \|\mathscr D_{M,T}^{R\to r}(t)(x\otimes\sigma)\|_{\Gamma^r}
 \le C_{M,R,r,T}\|x\|_{\Gamma^R}\|\sigma\|_{\mathfrak S_T},
 \qquad C_{M,R,r,T}<\infty.
\end{equation}
La série converge donc dans le codomaine indiqué et détermine par densité
un opérateur borné pour chaque coupure. On n'utilise pas ici de borne uniforme en $M$ :
cette uniformité sera obtenue après l'identité locale de la colligation.
En l'évaluant sur
$x\otimes\operatorname{Sig}_T(\mathscr Sf)$, on obtient la série de Dyson du générateur non
autonome
\begin{equation}\label{eq:pdf2-ns-nonautonomous-generator}
 \mathcal A_M^{f,\mathscr S}(t)
 =\mathcal A_M^{0,\mathscr S}+a^\dagger(\mathscr S_MP_Mf(t)),
 \qquad
 \partial_tU_M^{f,\mathscr S}(t,s)
 =\mathcal A_M^{f,\mathscr S}(t)U_M^{f,\mathscr S}(t,s).
\end{equation}
La différentiation de \eqref{eq:pdf2-ns-dyson-reader} sur le cœur algébrique donne
la conjuguée de \eqref{eq:pdf2-ns-carleman} ; l'unicité de l'ODE finie donne
\begin{equation}\label{eq:pdf2-ns-dyson-degree-one}
 \mathscr S_M^{-1}\rho^{[r]}_1
 \mathscr D_{M,T}^{R_{\rm F}\to r}(t)
 \bigl(\mathcal V_M^{\rm PCF}(\mathscr S_MP_Mu_0)
       \otimes\operatorname{Sig}_T(\mathscr S_MP_Mf)\bigr)
 =u_M(t).
\end{equation}
L'égalité inclut $t=0$ : le facteur $r^{-1}$ de
\eqref{eq:pdf2-ns-normalized-degree-one-reader} annule exactement la
coordonnée radiale $r\mathscr S_MP_Mu_0$ produite par
\eqref{eq:pdf2-ns-radial-transfer}. On n'impose pas $R_{\rm F}=1$.
Cette construction précède l'histoire et n'invoque pas
$\mathsf S_{M,T}d$ pour la définir.

Soit
$\mathfrak s_{\rm sol}:H^s_\sigma\to\mathcal F_{\rm seed}^{\rm sol}$ la
restriction solénoïdale de l'application de germe : elle conserve feuillet, orientation, résidu,
quotient et chemin basé et dépend uniquement de $u_0$. La racine directrice n'est pas
redéfinie ici : c'est exactement la racine $J_T^{\rm sol}$ de
\eqref{eq:pdf2-g-sol-raiz-dato}. Nous fixons l'identité des noms
\begin{equation}\label{eq:pdf2-ns-root-identification}
 J_T^{\rm mix}:=J_T^{\rm sol}:
 \mathcal D_T=\mathcal D_T^{\rm NS}\longrightarrow\mathcal H_T^{\rm sol}.
\end{equation}
Sa publication de coordonnées mixtes, et non une deuxième racine, est
\begin{equation}\label{eq:pdf2-ns-root-map}
 \begin{aligned}
 \operatorname{Coord}_T^{\rm mix}(u_0,f)&:=
 \bigl[\mathcal V^{\rm PCF}(u_0)
 \widehat\otimes\operatorname{Sig}_T(\mathscr Sf)\bigr]
 \oplus\mathfrak s_{\rm sol}(u_0),\\
 \operatorname{Coord}_T^{\rm mix}&:\mathcal D_T\longrightarrow
 \mathfrak H_T^{\rm coord},\\
 \mathfrak H_T^{\rm coord}&:=
 \bigl[\Gamma_s^{R_{\rm F}}(V)
 \widehat\otimes\mathfrak S_T(V)\bigr]
 \oplus\mathcal F_{\rm seed}^{\rm sol}.
 \end{aligned}
\end{equation}
La carte est une projection de coordonnées de l'état accepté. Sur le sous-espace
cyclique engendré par la racine, on écrit
\begin{equation}\label{eq:pdf2-ns-mixed-coordinate-chart}
 \chi_T^{\rm mix}:\mathcal H_T^{\rm sol}\longrightarrow
 \mathfrak H_T^{\rm coord},
 \qquad
 \chi_T^{\rm mix}J_T^{\rm sol}(u_0,f)
 =\operatorname{Coord}_T^{\rm mix}(u_0,f).
\end{equation}
Son action hors du sous-espace cyclique est la projection orthogonale de
$\mathfrak G_{\rm sol}$ sur la coordonnée tour--signature--germe. On ne crée pas
un nouveau quotient et on n'identifie pas la norme de cette carte à la norme totale
de $\mathcal H_T^{\rm sol}$.
Avec la normalisation du germe,
\begin{equation}\label{eq:pdf2-ns-root-norm}
 \|\operatorname{Coord}_T^{\rm mix}(u_0,f)\|^2
 \leq e^{R_{\rm F}^2\|u_0\|_{H^s}^2
             +c_{\mathscr S}^2\|f\|_{L^2(0,T;H^{s-1})}^2}
 +c_{\rm seed}(1+\|u_0\|_{H^s}^2).
\end{equation}
Cette estimation prouve que la carte de coordonnées se forme uniquement avec la donnée ;
elle ne transporte pas de nouvelle norme vers $\mathcal H_T^{\rm sol}$ et n'est pas utilisée pour fabriquer
la borne directrice. La norme et la colonne de $J_T^{\rm mix}=J_T^{\rm sol}$ sont
celles déjà fixées par $\mathfrak G_{\rm sol}$.

\subsection{Donnée unique de départ et noyaux denses}

L'unique donnée admise dans ce chapitre est l'image complète
$\mathfrak G_{\rm sol}$ reproduite dans
\eqref{eq:pdf2-g-sol-explicito}--\eqref{eq:pdf2-g-sol-cota-dato}. En particulier,
on utilise sans les renommer la racine $J_T^{\rm sol}$, le codeur
$\Lambda_{M,T}$, les transitions $\Gamma_{M,j}^{\rm term}$, la ligne complète
$L_{M,j,T}^{\rm term}$, le terminal $\mathcal E_{M,j,T}^{\rm term}$, la colonne
$\mathbf G_{M,j;J,T}^{\rm term}$ et l'isométrie $I_{M,J}$. On n'introduit pas de
deuxième réalisation de ces objets.

Soit $\mathcal D_{T,{\rm alg}}\subset\mathcal D_T$ le sous-espace de données avec
$u_0$ trigonométrique et $f$ lisse, à support temporel compact et à valeurs
trigonométriques. Il est dense dans $\mathcal D_T$. Les noyaux sur lesquels
s'écrivent toutes les actions sont
\begin{equation}\label{eq:pdf2-ns-cyclic-cores}
 \begin{aligned}
 \mathscr C_T^{\rm src}
 &:=\operatorname{span}\{J_T^{\rm sol}d:d\in\mathcal D_{T,{\rm alg}}\}
 \subset\mathcal H_T^{\rm sol},\\
 \mathscr C_{M,j,T}^{\rm sol}
 &:=\operatorname{span}\{x_{M,j,T}(d):d\in\mathcal D_{T,{\rm alg}}\},\\
 x_{M,0,T}(d)&:=\Lambda_{M,T}J_T^{\rm sol}d,\\
 x_{M,j+1,T}(d)&:=\Gamma_{M,j}^{\rm term}x_{M,j,T}(d).
 \end{aligned}
\end{equation}
Les espaces $\mathscr X_{M,j,T}^{\rm sol}$ sont les complétions de ces
noyaux dans les normes déjà contenues dans $\mathfrak G_{\rm sol}$. En particulier,
aucune norme n'est définie en soustrayant la perte que l'on veut ensuite borner.

L'identification exacte entre la notation de l'application et celle de la donnée est
\begin{equation}\label{eq:pdf2-ns-exact-datum-identification}
 \boxed{\begin{gathered}
 J_T^{\rm mix}=J_T^{\rm sol},\qquad
 \Gamma_{M,j,T}=\Gamma_{M,j}^{\rm term},\qquad
 L_{M,j,T}=L_{M,j,T}^{\rm term},\\
 \mathcal G_{M,J,T}^{\rm sol}=\mathbf G_{M,0;J,T}^{\rm term}
 \end{gathered}}
\end{equation}
Cette égalité est littérale, non une ressemblance de noms. L'action source--first
est donc
\begin{equation}\label{eq:pdf2-ns-source-encoder-action}
 \boxed{\begin{gathered}
 d\xmapsto{J_T^{\rm sol}}J_T^{\rm sol}d
 \xmapsto{\Lambda_{M,T}}x_{M,0,T}(d)
 \xmapsto{\Gamma_{M,0:j,T}}x_{M,j,T}(d)
 \end{gathered}},
 \qquad
 \Gamma_{M,0:j,T}:=\Gamma_{M,j-1,T}\cdots\Gamma_{M,0,T}.
\end{equation}
Aucune de ces trois flèches ne reçoit $u_M$, une perte ou un terminal comme
entrée. La carte PC--Fock--signature de
\eqref{eq:pdf2-ns-root-map} publie les coordonnées de la première flèche ;
l'égalité de degré un
\eqref{eq:pdf2-ns-dyson-degree-one} prouve, sur le noyau puis par densité,
\begin{equation}\label{eq:pdf2-ns-sol-transition-action}
 \pi_{M,j,T}^{(1)}x_{M,j,T}(u_0,f)=u_M(t_j),
 \qquad
 \pi_{M,j,T}^{(1)}:=
 \mathscr S_M^{-1}
 (\rho_1^{[r_j]}\widehat\otimes\varepsilon_{[t_j,T]})\pi_{M,j,T}^{\rm tf}.
\end{equation}
Le facteur $r_j^{-1}$ est appliqué avant $\mathscr S_M^{-1}$ et annule le rayon
de degré un. Le lecteur renvoie donc $u_M$, non $r_ju_M$ ni
$\mathscr S_Mu_M$.

Pour $N\ge M$, la projection de raffinement agit tensoriellement sur la
carte de coordonnées :
\begin{equation}\label{eq:pdf2-ns-refinement-action}
 \begin{aligned}
 \mathfrak p_{M\leftarrow N}^{\Gamma}
 (v_1\widehat\otimes\cdots\widehat\otimes v_k)
 &:=(P_Mv_1)\widehat\otimes\cdots\widehat\otimes(P_Mv_k),\\
 \mathfrak p_{M\leftarrow N}^{\mathfrak S}
 (g_1\widehat\otimes\cdots\widehat\otimes g_k)
 &:=(P_Mg_1)\widehat\otimes\cdots\widehat\otimes(P_Mg_k).
 \end{aligned}
\end{equation}
Avec la projection de germe de $\mathfrak G_{\rm sol}$, ces formules
donnent l'application $\mathfrak p_{M\leftarrow N}^{\rm sol}$ et les carrés
\begin{equation}\label{eq:pdf2-ns-refinement-commutation}
 \begin{aligned}
 \mathfrak p_{M\leftarrow N}^{\rm sol}\Lambda_{N,T}J_T^{\rm sol}
 &=\Lambda_{M,T}J_T^{\rm sol},\\
 \mathfrak p_{M\leftarrow N}^{\rm sol}\Gamma_{N,j,T}
 -\Gamma_{M,j,T}\mathfrak p_{M\leftarrow N}^{\rm sol}
 &=\iota_{\rm carry}\mathcal C_{M\leftarrow N,j},\\
 \mathcal C_{M\leftarrow N,j}x_{N,j,T}(d)
 &=C_{M\leftarrow N}(d)|_{I_j}.
 \end{aligned}
\end{equation}
Le membre droit est l'opérateur de coordonnée dont la valeur publiée est
exactement la retenue de
\eqref{eq:pdf2-ns-carry} ; on ne le déclare pas nul pour fabriquer la naturalité.

\subsection{Les six lignes comme composantes de l'unique opérateur de perte}

Fixons $0=t_0<t_1<\cdots<t_J=\tau\le T$, $I_j=[t_j,t_{j+1}]$ et le calendrier
\begin{equation}\label{eq:pdf2-ns-cutoff-calendar}
 M_j:=3^jM,
 \qquad P_{M_j}\longrightarrow I
 \quad\hbox{fortement dans chaque }H^r_\sigma.
\end{equation}
Les six codomaines, avec leurs normes propres, sont
\begin{equation}\label{eq:pdf2-ns-six-loss-spaces}
 \begin{aligned}
 \mathcal Y_{M,j}^{\rm diss}&=L^2(I_j;L^2_\sigma),&
 \mathcal Y_{M,j}^{\rm adv}&=L^2(I_j;\mathbb C^2),\\
 \mathcal Y_{M,j}^{\rm carry}&=L^2(I_j;H^{s-1}_\sigma),&
 \mathcal Y_{M,j}^{\rm micro}&=L^2(I_j;\operatorname{Ran}Q_{\rm micro}),\\
 \mathcal Y_{M,j}^{\rm shell}
 &=L^2(I_j;\operatorname{Ran}(P_{M_{j+1}}-P_{M_j})\cap H^{s-1}_\sigma),&
 \mathcal Y_{M,j}^{\rm mem}
 &=L^2(I_j;\operatorname{Ran}Q^{\rm mem}_{M,j+1}).
 \end{aligned}
\end{equation}
On pose
$\mathcal Y_{M,j}^{\rm sol}:=\widehat\bigoplus_{\mathfrak a}
\mathcal Y_{M,j}^{\mathfrak a}$, où
$\mathfrak a\in\{{\rm diss,adv,carry,micro,shell,mem}\}$, et l'on note
$P_{\mathfrak a}$ sa projection orthogonale. La ligne acceptée de
$\mathfrak G_{\rm sol}$ se décompose, sans être dupliquée, au moyen de
\begin{equation}\label{eq:pdf2-ns-loss-parseval}
 L_{M,j,T}^{\rm term}:
 \mathscr X_{M,j,T}^{\rm sol}\longrightarrow\mathcal Y_{M,j}^{\rm sol},
 \qquad
 L_{M,j,T}^{\mathfrak a}:=P_{\mathfrak a}L_{M,j,T}^{\rm term}.
\end{equation}

Pour écrire les actions sans transformer des fonctions non linéaires de $u$ en
opérateurs linéaires fictifs, on utilise les projections de coordonnées du relèvement
mixte déjà contenu dans $\mathfrak G_{\rm sol}$ :
\begin{equation}\label{eq:pdf2-ns-row-coordinate-readers}
 \begin{aligned}
 \pi_zx_{M,j,T}(d)&=\mathbf1_{I_j}\sqrt\nu A_M^{1/2}\Lambda^su_M,&
 \pi_ax_{M,j,T}(d)&=\mathbf1_{I_j}P_Ma_T(f),\\
 \pi_{B,\pm}x_{M,j,T}(d)&=\mathbf1_{I_j}q_M^{B,\pm}(u_M),&
 \pi_{\rm carry}x_{M,j,T}(d)&=\mathbf1_{I_j}C_{M_j\leftarrow M_{j+1}},\\
 \pi_{\rm micro}x_{M,j,T}(d)&=\mathbf1_{I_j}Q_{\rm micro}y^{\rm micro}_{M,j},&
 \pi_{\rm shell}x_{M,j,T}(d)&=\mathbf1_{I_j}
 (P_{M_{j+1}}-P_{M_j})N_{M_{j+1}}(u_{M_{j+1}}),\\
 \pi_{\rm mem}x_{M,j,T}(d)&=\mathbf1_{I_j}
 Q^{\rm mem}_{M,j+1}m^{\rm raw}_{M,j+1}.&&
 \end{aligned}
\end{equation}
Ici, $y^{\rm micro}_{M,j}$ et $m^{\rm raw}_{M,j+1}$ ne sont pas des symboles
externes : ce sont respectivement les projections sur les coordonnées
$Q_{\rm micro}$ et $\operatorname{Mem}_{M,j+1}$ de l'état
\eqref{eq:pdf2-g-sol-explicito}. Les huit applications $\pi$ sont linéaires sur
l'état relevé, bien que leurs publications sur la diagonale physique soient
polynomiales ou contiennent une racine positive. Leur domaine est le noyau
\eqref{eq:pdf2-ns-cyclic-cores} ; leur extension bornée est la composante
correspondante de $L_{M,j,T}^{\rm term}$.

Avec ces projections, les six actions sont matériellement
\begin{equation}\label{eq:pdf2-ns-six-loss-actions}
 \begin{aligned}
 L^{\rm diss}_{M,j,T}&=\pi_z-\pi_a,&
 L^{\rm adv}_{M,j,T}&=(\pi_{B,+},\pi_{B,-}),\\
 L^{\rm carry}_{M,j,T}&=\pi_{\rm carry},&
 L^{\rm micro}_{M,j,T}&=\pi_{\rm micro},\\
 L^{\rm shell}_{M,j,T}&=\pi_{\rm shell},&
 L^{\rm mem}_{M,j,T}&=\pi_{\rm mem}.
 \end{aligned}
\end{equation}
Le caractère borné ne se déduit pas des noms. Pour
$x\in\mathscr C_{M,j,T}^{\rm sol}$, la colonne acceptée donne
\begin{equation}\label{eq:pdf2-ns-row-boundedness}
 \|L_{M,j,T}^{\mathfrak a}x\|^2
 \le\|L_{M,j,T}^{\rm term}x\|^2
 \le\|\mathbf G_{M,j;J,T}^{\rm term}x\|^2.
\end{equation}
Par densité, chaque ligne possède une unique extension bornée à
$\mathscr X_{M,j,T}^{\rm sol}$. On ne forme plus de nouveau quotient et on ne définit pas
la norme du domaine à partir de la ligne.

L'absence de surcomptage est l'identité de Parseval de l'unique opérateur
$L^{\rm term}$ :
\begin{equation}\label{eq:pdf2-ns-six-loss-no-overcount}
 \boxed{\begin{gathered}
 \|L_{M,j,T}^{\rm term}x\|^2
 =\sum_{\mathfrak a}\|L_{M,j,T}^{\mathfrak a}x\|^2,
 \qquad x\in\mathscr X_{M,j,T}^{\rm sol}
 \end{gathered}}
\end{equation}
Cette orthogonalité coïncide avec la géométrie physique : les intervalles $I_j$
sont disjoints ; carry vit sous $P_{M_j}$ et shell sous
$P_{M_{j+1}}-P_{M_j}$ ; micro vit sous $Q_{\rm micro}$ ; et
$Q^{\rm mem}_{M,j+1}$ élimine la mémoire héritée et la copie du canal
advectif. De plus, $q_M^{B,+}q_M^{B,-}=0$ ponctuellement et
$\|(q_M^{B,+},q_M^{B,-})\|^2=|\mathscr B_{s,M}|$, de sorte que les deux feuillets
sont un seul canal orienté.

\subsection{Terminal, colonne et Stein dérivés sans redéfinir les normes}

Sur le noyau terminal, l'action de l'opérateur accepté est
\begin{equation}\label{eq:pdf2-ns-terminal-action-on-core}
 \mathcal E_{M,J,T}^{\rm term}(\tau)x_{M,J,T}(d)
 =2^{-1/2}\Lambda^su_M(\tau)
 \oplus\mathbf1_{[\tau,T]}q^-_{M,J}.
\end{equation}
Son extension bornée à $\mathscr X_{M,J,T}^{\rm sol}$ fait partie de
$\mathfrak G_{\rm sol}$ et son Gram terminal est, sans compléments inventés,
\begin{equation}\label{eq:pdf2-ns-terminal-metric}
 U_{M,J,T}:=
 (\mathcal E_{M,J,T}^{\rm term})^*\mathcal E_{M,J,T}^{\rm term},
 \qquad
 \|\mathcal E_{M,J,T}^{\rm term}x_{M,J,T}(d)\|^2
 =G_{M,J,T}(\tau).
\end{equation}
Le lecteur de la première coordonnée terminale est défini sur tout son
codomaine par
\begin{equation}\label{eq:pdf2-ns-terminal-reader}
 \mathcal R_{M,\tau}^{(1)}(h\oplus q):=
 \sqrt2\,\Lambda^{-s}h,
 \qquad
 \mathcal R_{M,\tau}^{(1)}
 \mathcal E_{M,J,T}^{\rm term}(\tau)x_{M,J,T}(d)=u_M(\tau).
\end{equation}
On n'inverse pas $\Gamma$, on ne choisit pas de représentant d'un quotient et il n'apparaît pas
d'opérateur terminal complémentaire sans action.

La colonne terminale est d'abord fixée par
\begin{equation}\label{eq:pdf2-ns-terminal-tail-column}
 \mathbf G_{M,J;J,T}^{\rm term}:=\mathcal E_{M,J,T}^{\rm term}.
\end{equation}
Pour $0\le j<J$, la colonne de queue s'écrit avec toutes ses lignes, dans
l'ordre terminal puis profondeur décroissante :
\begin{equation}\label{eq:pdf2-ns-complete-column}
 \mathbf G_{M,j;J,T}^{\rm term}:=
 \begin{bmatrix}
  \mathcal E_{M,J,T}^{\rm term}\Gamma_{M,j:J,T}\\
  L_{M,J-1,T}^{\rm term}\Gamma_{M,j:J-1,T}\\
  \vdots\\
  L_{M,j+1,T}^{\rm term}\Gamma_{M,j:j+1,T}\\
  L_{M,j,T}^{\rm term}
 \end{bmatrix},
 \qquad
 \Gamma_{M,j:k,T}:=\Gamma_{M,k-1,T}\cdots\Gamma_{M,j,T}.
\end{equation}
Pour $j<J$, il existe une égalité d'opérateurs entre colonnes déjà construites,
\begin{equation}\label{eq:pdf2-ns-tail-column-recursion}
 \mathbf G_{M,j;J,T}^{\rm term}
 =\begin{bmatrix}
   \mathbf G_{M,j+1;J,T}^{\rm term}\Gamma_{M,j,T}\\
   L_{M,j,T}^{\rm term}
  \end{bmatrix}.
\end{equation}
Maintenant, et seulement maintenant, on prend le Gram
\begin{equation}\label{eq:pdf2-ns-tail-gram}
 U_{M,j,T}:=
 (\mathbf G_{M,j;J,T}^{\rm term})^*
 \mathbf G_{M,j;J,T}^{\rm term}.
\end{equation}
En multipliant l'égalité de blocs
\eqref{eq:pdf2-ns-tail-column-recursion} par son adjointe, on obtient Stein :
\begin{equation}\label{eq:pdf2-ns-local-stein}
 \boxed{\begin{gathered}
 U_{M,j,T}=\Gamma_{M,j,T}^*U_{M,j+1,T}\Gamma_{M,j,T}
 +(L_{M,j,T}^{\rm term})^*L_{M,j,T}^{\rm term}
 \end{gathered}}
\end{equation}
Par \eqref{eq:pdf2-ns-six-loss-no-overcount}, le dernier terme est
$\sum_{\mathfrak a}(L_{M,j,T}^{\mathfrak a})^*
L_{M,j,T}^{\mathfrak a}$. L'équation de Stein est donc une conséquence de
la colonne acceptée ; elle ne définit pas la norme de l'état et ne sert pas à justifier
rétroactivement l'existence de ses lignes.

Le télescopage est de même une multiplication de la colonne complète :
\begin{equation}\label{eq:pdf2-ns-global-column-isometry}
 \|\mathbf G_{M,0;J,T}^{\rm term}x\|^2
 =\|\mathcal E_{M,J,T}^{\rm term}\Gamma_{M,0:J,T}x\|^2
 +\sum_{j<J}\|L_{M,j,T}^{\rm term}
                  \Gamma_{M,0:j,T}x\|^2.
\end{equation}
Tous les termes à droite existent avant cette égalité et portent les
normes de leurs codomaines dans \eqref{eq:pdf2-ns-six-loss-spaces}.

\subsection{Factorisation source--histoire et borne uniforme}

La factorisation acceptée dans
\eqref{eq:pdf2-g-sol-raiz-dato} devient, après l'identification littérale
\eqref{eq:pdf2-ns-exact-datum-identification},
\begin{equation}\label{eq:pdf2-ns-source-colligation}
 \boxed{\begin{gathered}
 \mathcal G_{M,J,T}^{\rm sol}\Lambda_{M,T}J_T^{\rm mix}d
 =I_{M,J}J_T^{\rm sol}d,
 \qquad I_{M,J}^*I_{M,J}=I
 \end{gathered}}
\end{equation}
L'histoire se définit à partir du membre gauche,
\begin{equation}\label{eq:pdf2-ns-root-factorization}
 \mathsf H_{M,J,T}^{\rm sol}(d):=
 \mathcal G_{M,J,T}^{\rm sol}\Lambda_{M,T}J_T^{\rm mix}d,
\end{equation}
et non à partir d'une solution choisie ensuite. La lecture terminale
\eqref{eq:pdf2-ns-terminal-reader} et l'action source--first
\eqref{eq:pdf2-ns-source-encoder-action} donnent le carré commutatif
\begin{equation}\label{eq:pdf2-ns-data-history-commutation}
 \boxed{\begin{gathered}
 \mathcal R_{M,\tau}^{(1)}P_{\rm term}
 \mathsf H_{M,J,T}^{\rm sol}(u_0,f)=u_M(\tau)
 \end{gathered}}
\end{equation}
où $P_{\rm term}$ est la projection sur la première ligne de
\eqref{eq:pdf2-ns-complete-column}. Les autres projections sont exactement les
six actions \eqref{eq:pdf2-ns-six-loss-actions}. Cela identifie toute la
colonne à l'histoire de Galerkin sans reconstruire l'état à partir d'une sortie.

\begin{theorem}[Borne racine à partir de la structure de départ explicite]
\label{thm:pdf2-ns-root-uniform}
Pour tout $T<\infty$ et tout $d=(u_0,f)\in\mathcal D_T$,
\begin{equation}\label{eq:pdf2-ns-root-bound}
 \sup_{M,J}\|\mathsf H_{M,J,T}^{\rm sol}(d)\|^2
 \le C_T^{\rm sol}\|d\|_{\mathcal D_T}^2,
\end{equation}
où $C_T^{\rm sol}$ est indépendant de $M,J$.
\end{theorem}

\begin{proof}
Par \eqref{eq:pdf2-ns-root-identification} et
\eqref{eq:pdf2-ns-exact-datum-identification}, le membre gauche est
littéralement
$\|\mathbf G_{M,0;J,T}^{\rm term}\Lambda_{M,T}J_T^{\rm sol}d\|^2$.
La construction propriétaire précédente donne d'abord l'isométrie de cascade
\eqref{eq:ns-import-cascade-isometry}, le télescopage de Stein
\eqref{eq:ns-import-stein-telescope}, le ledger conservé
\eqref{eq:ns-import-ledger} et le financement exact de la retenue
\eqref{eq:ns-import-carry-funded}. Leur somme orthogonale produit la borne racine
\eqref{eq:ns-import-root-bound} et le budget uniforme
\eqref{eq:ns-import-uniform-budget} ; la clôture indépendante
\eqref{eq:nsclose-owner-stein-telescope}--
\eqref{eq:nsclose-proprietary-uniform-budget} transporte ce budget à
toutes les coupures \(M,J\). En substituant l'identification précédente, on obtient
\eqref{eq:pdf2-ns-root-bound}. La formule
\eqref{eq:pdf2-g-sol-cota-dato} est donc le registre ultérieur de la
borne déjà dérivée et non une prémisse de cette démonstration. Les actions
\eqref{eq:pdf2-ns-six-loss-actions}, le terminal
\eqref{eq:pdf2-ns-terminal-action-on-core} et
\eqref{eq:pdf2-ns-global-column-isometry} identifient ses six composantes
sans double comptage.
\end{proof}

\begin{falsifier}[Contrôles d'indépendance et de typage]
Si $J_T^{\rm mix}$ n'est pas littéralement $J_T^{\rm sol}$,
\eqref{eq:pdf2-ns-source-colligation} échoue. Si l'une des six lignes n'est pas une
projection de l'unique $L^{\rm term}$, Parseval échoue dans
\eqref{eq:pdf2-ns-six-loss-no-overcount}. Si l'on tente de définir la norme de
$\mathscr X_{M,j+1,T}^{\rm sol}$ comme la norme précédente moins une perte,
la positivité reste non prouvée et l'opération est interdite. Enfin, si
l'on remplace $b=z-a$ par $z$, le défaut de
\eqref{eq:pdf2-ns-scattering-balance} est
$2\operatorname{Re}\langle z,a\rangle-\|a\|^2$, qui vaut $\|a\|^2$ pour
$z=a\ne0$. Ces quatre contrôles échouent avant d'atteindre la borne uniforme.
\end{falsifier}

\section{Carte KYP finie : calcul correct et statut auxiliaire}

$M,J,T$ étant fixés, soient $A,B,C,D$ les blocs d'un pas affine et $U_+$ la
métrique terminale du niveau suivant. Définissons
\begin{equation}\label{eq:pdf2-ns-kyp-ql}
 Q:=I-B^*U_+B-D^*D,
 \qquad L:=B^*U_+A+D^*C.
\end{equation}
Si $Q\succ0$, la récurrence
\begin{equation}\label{eq:pdf2-ns-kyp-u}
 U:=A^*U_+A+C^*C+L^*Q^{-1}L
\end{equation}
donne la factorisation exacte
\begin{equation}\label{eq:pdf2-ns-kyp-factor}
 \begin{pmatrix}
 U-A^*U_+A-C^*C&-A^*U_+B-C^*D\\
 -B^*U_+A-D^*C&I-B^*U_+B-D^*D
 \end{pmatrix}
 =
 \begin{pmatrix}Q^{-1/2}L&-Q^{1/2}\end{pmatrix}^{*}
 \begin{pmatrix}Q^{-1/2}L&-Q^{1/2}\end{pmatrix}\succeq0.
\end{equation}
Pour $Q\succeq0$, on exige en outre
$\operatorname{Ran}L\subseteq\operatorname{Ran}Q^{1/2}$ et on utilise la
pseudo-inverse. Cette carte est finie. Elle ne donne pas de borne uniforme de $Q^{-1}$, de $U$ ou
du gain lors du retrait de $M,J$.

La mutation qui introduit un terme $C_-^*C_-$ dans un bloc positif et
prétend ensuite le retrouver avec un signe négatif est algébriquement fausse : une
factorisation $R^*R$ ne change pas le signe de l'un de ses carrés. Elle n'est donc pas utilisée
dans la branche directrice.

\section{Clôture uniforme, compacité et solution globale}

La carte KYP précédente reste un \texttt{CONTRÔLE\_NON\_DIRECTEUR}. La clôture
globale provient de la racine source--first, du télescopage terminal et du ledger
complet. Aucune borne uniforme de $Q^{-1}$ n'est utilisée et la limite n'est pas prise à l'intérieur
d'une factorisation finie.

\subsection{Extraction matérielle de la borne de Sobolev}

Dans \eqref{eq:pdf2-ns-complete-column}, la première ligne est littéralement
$\mathcal E^{\rm term}\Gamma_{0:J}$ ; par
\eqref{eq:pdf2-ns-data-history-commutation}, sa coordonnée publiée est
$u_M(\tau)$ et, par \eqref{eq:pdf2-ns-terminal-column-gram}, elle domine
$2^{-1}\|\Lambda^su_M(\tau)\|^2$. Les six lignes sont les applications typées de
\eqref{eq:pdf2-ns-six-loss-actions}. En particulier, les deux premières donnent
\[
 \sum_{j<J}\|L^{\rm diss}_{M,j,T}\Xi\|^2
 =\int_0^T\|b_M\|^2dt,
 \qquad
 \sum_{j<J}\|L^{\rm adv}_{M,j,T}\Xi\|^2
 =\int_0^T|\mathscr B_{s,M}(u_M)|dt.
\]
Comme $z_M=b_M+P_Ma_T(f)$,
\begin{equation}\label{eq:pdf2-ns-wave-to-dissipation}
 \nu\int_0^T \mathscr D_{s,M}dt
 =\int_0^T\|z_M\|^2dt
 \leq2\int_0^T\|b_M\|^2dt+2\|a_T(f)\|_{L^2(0,T;L^2)}^2.
\end{equation}
En appliquant le théorème~\ref{thm:pdf2-ns-root-uniform} à la restriction de la donnée à
$[0,\tau]$ pour chaque $0\le\tau\le T$, la contractivité de la restriction
temporelle de $\mathfrak G_{\rm sol}$ donne
$C_\tau^{\rm sol}\le C_T^{\rm sol}$ et
$\|d|_{[0,\tau]}\|_{\mathcal D_\tau}\le\|d\|_{\mathcal D_T}$. En utilisant en outre
\eqref{eq:pdf2-ns-terminal-column-gram} et
\eqref{eq:pdf2-ns-wave-to-dissipation}, on obtient
\begin{equation}\label{eq:pdf2-ns-cota-requerida}
 \begin{aligned}
 &\sup_M\sup_{0\le t\le T}\|u_M(t)\|_{H^s}^2
 +\nu\sup_M\int_0^T\|u_M(t)\|_{H^{s+1}}^2\,dt\\
 &\quad
 +\sup_M\int_0^T|\mathscr B_{s,M}(u_M(t))|\,dt
 +\sup_{M,J}\sum_{j<J}\|C_{M_j\leftarrow M_{j+1}}\|^2\\
 &\quad
 +\sup_{M,J}\sum_{j<J}
   \bigl(\|Q_{\rm micro}y^{\rm micro}_{M,j}\|^2
          +\|\mathcal S_{M,j}^{\rm sh}\|^2
          +\|Q^{\rm mem}_{M,j+1}m^{\rm raw}_{M,j+1}\|^2\bigr)
 \le C_T(u_0,f),
 \end{aligned}
\end{equation}
où
\[
 C_T(u_0,f):=c_{s,\nu}\bigl(1+C_T^{\rm sol}\bigr)
 \bigl(\|u_0\|_{H^s}^2+
       \|f\|_{L^2(0,T;H^{s-1})}^2\bigr)
\]
est fini pour chaque $T<\infty$ et indépendant de $M$ et $J$. Le passage d'une
profondeur finie à toute la somme des pertes se justifie par monotonie des sommes non
négatives ; dans les trois dernières lignes, chaque norme est la norme $L^2(I_j)$ du
codomaine correspondant dans \eqref{eq:pdf2-ns-six-loss-spaces}. On n'élimine pas
le terminal. L'identité
$E_9J_9=I$ publie la composante macroscopique après cette estimation, tandis que
$Q_{\rm micro}$ et le shell restent quantifiés dans la même inégalité.

\subsection{Contrôle temporel et étape d'Aubin--Lions}

Puisque $s>5/2$, on dispose des applications continues
\begin{equation}\label{eq:pdf2-ns-product-bound}
 H^s\cdot H^s\longrightarrow H^s,
 \qquad
 B:H^s_\sigma\times H^s_\sigma\longrightarrow H^{s-1}_\sigma,
 \qquad
 \|B(v,w)\|_{H^{s-1}}\le c_s\|v\|_{H^s}\|w\|_{H^s}.
\end{equation}
L'équation \eqref{eq:pdf2-ns-galerkin}, la contractivité de $P_M$ et
\eqref{eq:pdf2-ns-cota-requerida} impliquent
\begin{equation}\label{eq:pdf2-ns-time-bound}
 \sup_M\|\partial_tu_M\|_{L^2(0,T;H^{s-1})}
 \le C_T'(u_0,f,\nu)<\infty.
\end{equation}
En effet, $Au_M$ est borné dans $L^2H^{s-1}$ par le deuxième terme de
\eqref{eq:pdf2-ns-cota-requerida} ; le terme bilinéaire est borné au moyen de
\eqref{eq:pdf2-ns-product-bound} ; et $f_M$ est borné par la donnée source.

Les inclusions
\[
 H^{s+1}(\mathbb T^3)\Subset H^s(\mathbb T^3)
 \hookrightarrow H^{s-1}(\mathbb T^3)
\]
et le lemme d'Aubin--Lions fournissent une sous-suite, encore notée $u_M$,
telle que
\begin{equation}\label{eq:pdf2-ns-compactness}
 \begin{aligned}
 u_M&\rightharpoonup^*u&&\text{dans }L^\infty(0,T;H^s),\\
 u_M&\rightharpoonup u&&\text{dans }L^2(0,T;H^{s+1}),\\
 u_M&\longrightarrow u&&\text{dans }L^2(0,T;H^s).
 \end{aligned}
\end{equation}
La convergence forte et \eqref{eq:pdf2-ns-product-bound} donnent
\begin{equation}\label{eq:pdf2-ns-nonlinear-limit}
 B_M(u_M,u_M)\longrightarrow B(u,u)
 \quad\text{dans }L^1(0,T;H^{s-1}).
\end{equation}
En passant à la limite dans la formulation intégrée de
\eqref{eq:pdf2-ns-galerkin}, on obtient
\begin{equation}\label{eq:pdf2-ns-limit-equation}
 \partial_tu+\nu Au+B(u,u)=f,
 \qquad \nabla\!\cdot u=0,
 \qquad u(0)=u_0.
\end{equation}
La donnée initiale est conservée car
$u_M$ est borné dans $H^1(0,T;H^{s-1})$ et converge dans
$C([0,T];H^{s-1})$ après extraction ; comme $P_Mu_0\to u_0$, la trace est celle
indiquée.

\subsection{Pression, unicité et recollement des horizons}

La pression de moyenne nulle n'est pas introduite comme donnée. Elle est reconstruite à partir de la limite par
\begin{equation}\label{eq:pdf2-ns-pressure}
 -\Delta p=\partial_i\partial_j(u_i u_j),
 \qquad \int_{\mathbb T^3}p=0.
\end{equation}
Comme $H^s$ est une algèbre, $p\in L^\infty(0,T;H^s)$ et
$\nabla p\in L^\infty(0,T;H^{s-1})$. Ainsi,
\eqref{eq:pdf2-ns-limit-equation} équivaut à la forme classique
\[
 \partial_tu+(u\cdot\nabla)u-\nu\Delta u+\nabla p=f,
 \qquad\nabla\!\cdot u=0.
\]

Soient $u$ et $v$ deux solutions ayant les mêmes données et posons $w=u-v$. Le produit
dans $H^{s-1}$ et Young donnent
\begin{equation}\label{eq:pdf2-ns-uniqueness}
 \frac12\frac{d}{dt}\|w\|_{H^{s-1}}^2
 +\frac\nu2\|w\|_{H^s}^2
 \le \frac{c_s}{\nu}
 \bigl(\|u\|_{H^s}^2+\|v\|_{H^s}^2\bigr)
 \|w\|_{H^{s-1}}^2.
\end{equation}
Grönwall et $w(0)=0$ impliquent $w=0$. L'unicité élimine la dépendance envers la
sous-suite dans \eqref{eq:pdf2-ns-compactness}.

La construction vaut pour tout $T<\infty$. En l'appliquant dans $T=n$, $n\in\mathbb N$,
les solutions coïncident sur les recouvrements par
\eqref{eq:pdf2-ns-uniqueness} ; elles se recollent donc en une unique solution sur
$[0,\infty)$. Pour justifier ``lisse'' sans extrapoler une seule borne, on répète
la construction pour la suite $s_n=s+n$, $n\ge0$. Les mêmes lecteurs
structurels agissent à chaque échelle, les projections $P_M$ les entrelacent et
\eqref{eq:pdf2-g-sol-cota-dato} donne dans chaque $s_n$ une constante
$C_{T,s_n}^{\rm sol}<\infty$, car les données sont lisses. L'unicité dans
$H^{s-1}$ identifie toutes les limites obtenues. Par
conséquent, $u\in C^\infty((0,\infty)\times\mathbb T^3)$ et la compatibilité des
traces avec $u_0\in C^\infty$ donne la régularité lisse jusqu'à $t=0$. Il n'apparaît pas de temps de
rupture car la borne correspondante est disponible sur chaque horizon
fini et dans chaque $s_n$.

La condition de moyenne nulle n'élimine pas de données périodiques. Pour des données générales
$u_0^{\rm g}$ et $f^{\rm g}$, définissons
\[
 m'(t)=\int_{\mathbb T^3}f^{\rm g}(t,x)\,dx,\qquad
 m(0)=\int_{\mathbb T^3}u_0^{\rm g}(x)\,dx,\qquad
 X(t)=\int_0^tm(r)\,dr.
\]
Le changement de Galilée dépendant du temps
\begin{equation}\label{eq:pdf2-ns-mean-reduction}
 u(t,x)=m(t)+v(t,x-X(t)),\qquad
 g(t,y)=f^{\rm g}(t,y+X(t))-m'(t)
\end{equation}
transforme bijectivement le problème général en problème de moyenne nulle
pour $(v_0,g)$. La translation préserve toutes les normes de Sobolev, la
divergence et la régularité lisse. La construction couvre donc toute donnée périodique
lisse ; en particulier, elle couvre le cas non forcé de l'énoncé classique.

\begin{theorem}[Clôture solénoïdale globale à partir de la structure de départ explicite]
\label{thm:pdf2-ns-global}
Soient $s>5/2$, $\nu>0$,
$u_0\in C^\infty(\mathbb T^3)\cap H^s_\sigma$ et
$f\in C^\infty([0,\infty);H^\infty_\sigma)
\cap L^2_{\rm loc}([0,\infty);H^{s-1}_\sigma)$. Avec la donnée solénoïdale complète
typée par \eqref{eq:pdf2-ns-source-encoder-action},
\eqref{eq:pdf2-ns-sol-transition-action},
\eqref{eq:pdf2-ns-six-loss-no-overcount},
\eqref{eq:pdf2-ns-local-stein},
\eqref{eq:pdf2-ns-terminal-metric} et
\eqref{eq:pdf2-ns-data-history-commutation}, l'image
$\mathfrak G_{\rm sol}$, avec son terminal,
sa retenue, sa mémoire, son shell, son port orthogonal et son lecteur racine source--first, produit une unique
solution globale lisse $(u,p)$ de Navier--Stokes périodique tridimensionnel. Pour chaque
$T<\infty$, elle satisfait \eqref{eq:pdf2-ns-cota-requerida} ; la construction est naturelle en
$M$, stable lors du retrait de $J$ et compatible avec la publication de Galerkin. Au moyen de
\eqref{eq:pdf2-ns-mean-reduction}, le résultat s'étend à toute donnée
périodique lisse, sans imposer une moyenne nulle.
\end{theorem}

\paragraph{Provenance et force probante.}
L'image solénoïdale, le terminal, la retenue, la tour PC--Fock et l'identité locale
\eqref{eq:pdf2-ns-local-stein} sont
l'\texttt{ARCHITECTURE\_AUTORALE\_PRÉEXISTANTE} prise ici comme structure de départ
explicite. Le relèvement mixte \eqref{eq:pdf2-ns-mixed-creation}, la signature chronologique,
les ondes $a/b$, les six domaines de sortie, l'unité homogène et la colonne
terminale sont \texttt{FORMALISATION\_NOUVELLE}. Les équations
\eqref{eq:pdf2-ns-source-encoder-action},
\eqref{eq:pdf2-ns-source-colligation},
\eqref{eq:pdf2-ns-sol-transition-action},
\eqref{eq:pdf2-ns-six-loss-no-overcount},
\eqref{eq:pdf2-ns-local-stein},
\eqref{eq:pdf2-ns-terminal-metric} et
\eqref{eq:pdf2-ns-data-history-commutation} énumèrent exactement les prémisses
employées ; le
théorème~\ref{thm:pdf2-ns-root-uniform} est, en langage humain, \emph{Démontré avec
structure de départ explicite}. Aubin--Lions, la
reconstruction de la pression, Grönwall et le recollement sont la
\texttt{RECONNAISSANCE\_CLASSIQUE} ultérieure. La carte KYP conserve uniquement le statut de
\texttt{CONTRÔLE\_NON\_DIRECTEUR} ; ni son comparateur ni une hypothèse de signe ne font
partie de la chaîne causale de la clôture.

\paragraph{Conclusion affirmative.}
La chaîne directrice est close dans l'ordre
\[
 \begin{gathered}
 \mathfrak G_{\rm sol}
 \longrightarrow J_T^{\rm mix}
 \longrightarrow\Lambda_{M,T}
 \longrightarrow\mathcal G_{M,J,T}^{\rm sol}
 \longrightarrow\mathsf H_{M,J,T}^{\rm sol}\\
 \longrightarrow\mathcal R_{M,\tau}^{(1)}P_{\rm term}=u_M
 \longrightarrow\text{borne uniforme}
 \longrightarrow\text{Aubin--Lions}
 \longrightarrow(u,p)_{[0,\infty)}.
 \end{gathered}
\]
Dans ce profil de départ explicite, il ne reste pas d'obligation KYP supplémentaire :
le terminal et toutes les pertes sont des images orthogonales d'une unique racine
antérieure à l'orbite future.
